% !TEX TS-program = pdflatex
% !TeX program = pdflatex
% !TEX encoding = UTF-8
% !TEX spellcheck = fr

\documentclass[xcolor=table]{beamer}


%\usepackage{fullpage}
%\usepackage[left=2.8cm,right=2.2cm,top=2 cm,bottom=2 cm]{geometry}
\setbeamersize{text margin left=10pt,text margin right=10pt}


\usepackage{amsmath,amssymb} 
\usepackage{fontenc}
\usepackage{textcomp}
\usepackage[utf8]{inputenc}
\usepackage{CJKutf8}
\usepackage[french]{babel}
\usepackage{arabtex,acolor}
\usepackage{txfonts}
\usepackage{tipa}
\usepackage[]{graphicx}
\usepackage{multirow}
\usepackage{hyperref}
\usepackage{colortbl}
\usepackage{tabularray}
\usepackage{listingsutf8}
\usepackage{wrapfig}
\usepackage{multicol}
\usepackage[export]{adjustbox} %for images in table, also for frame
\usepackage[many]{tcolorbox}
\usepackage{wasysym}
\usepackage[french,lined]{algorithm2e}
\usepackage{alltt} %verbatim with commands
\usepackage{longtable}
\usepackage{tabu}


\definecolor{lightblue}{HTML}{D0D2FF}
\definecolor{lightyellow}{HTML}{FFFFAA}
\definecolor{darkblue}{HTML}{0000BB}
\definecolor{olivegreen}{HTML}{006600}
\definecolor{darkgreen}{HTML}{008B45} %009B55
\definecolor{violet}{HTML}{6600CC}
%\definecolor{deeppink}{HTML}{FF1493}
\definecolor{orangey}{HTML}{FFBB00}

\usetheme{Warsaw} % Antibes Boadilla Warsaw

\beamertemplatenavigationsymbolsempty


\let\oldcite\cite
\renewcommand{\cite}[1]{{\bfseries\color{orangey}\oldcite{#1}}}

\tcbuselibrary{listings}

\newcommand{\kurl}[1]{{\scriptsize\bfseries\color{orangey}\url{#1}}}


%\renewcommand{\baselinestretch}{1.5}

\def\supit#1{\raisebox{0.8ex}{\small\it #1}\hspace{0.05em}}

\AtBeginSection{%
	\begin{frame}
		\sectionpage
	\end{frame}
}

\newcommand{\rottext}[2]{%
	\rotatebox{90}{%
	\begin{minipage}{#1}%
		\raggedleft#2%
	\end{minipage}%
	}%
}


\institute{ %
Laboratoire de la Communication dans les Systèmes Informatiques (LCSI)

\vspace{6pt}
École  nationale Supérieure d'Informatique (ESI, ex. INI), Alger, Algérie
}
\author[ \textbf{\footnotesize\insertframenumber/\inserttotalframenumber} \hspace*{1cm} ESI - ARIES Abdelkrime (Master 2022/2023)] %
{ARIES Abdelkrime}
%\titlegraphic{\includegraphics[height=1cm]{../img/esi-logo.png}%\hspace*{4.75cm}~
	
	
	\date{Master 2022/2023} %\today

\titlegraphic{%
	\includegraphics[height=1cm]{../img/esi-logo.png}%
	\hspace{2cm}
	\includegraphics[height=1cm]{../img/lcsi-logo.png}%
}


%\setbeamertemplate{headline}{}



\newcommand{\keyword}[1]{\textcolor{red}{\bfseries\itshape #1}}
\newcommand{\expword}[1]{\textcolor{olivegreen}{#1}}
\newcommand{\optword}[1]{\textcolor{violet}{\bfseries #1}}

\makeatletter
\newcommand\mysphere{%
	\parbox[t]{10pt}{\raisebox{0.2pt}{\beamer@usesphere{item projected}{bigsphere}}}}
\makeatother

%\let\oldtabular\tabular
%\let\endoldtabular\endtabular
%\renewenvironment{tabular}{\rowcolors{2}{white}{lightblue}\oldtabular\rowcolor{blue}}{\endoldtabular}


\NoAutoSpacing %french autospacing after ":"

\def\graphpath{}

\newcommand{\changegraphpath}[1]{\def\graphpath{#1}}


\newcommand{\vgraphpage}[2][.84\textheight]{%
%	\begin{center}%
		\includegraphics[height=#1]{\graphpath #2}%
%	\end{center}%
}

\newcommand{\hgraphpage}[2][\textwidth]{%
%	\begin{center}%
		\includegraphics[width=#1]{\graphpath #2}%
%	\end{center}%
}

\newcommand{\graphpage}[2][]{%
	\includegraphics[#1]{\graphpath #2}%
}

\bibliographystyle{apalike}

\newcommand{\insertbibliography}[2]{
	\appendix
	\section*{Bibliographie}
	\nocite{#2}
%	\makeatletter % to change template
%	\setbeamertemplate{headline}[default] % not mandatory, but I though it was better to set it blank
%	\def\beamer@entrycode{\vspace*{-\headheight}} % here is the part we are interested in :)
%	\makeatother
	\begin{multicols*}{2}[\frametitle{\insertsection} \usebeamertemplate*{frametitle}]%\usebeamertemplate*{frametitle}\frametitle{Références}
		\tiny
		\bibliography{#1}
	\end{multicols*}
}

\definecolor{my-grey}{RGB}{233, 233, 233}

\newcommand{\insertlicence}{
	\begin{frame}[plain]
	\frametitle{Licence : CC-BY 4.0}
%	\framesubtitle{Licence: CC-BY-NC 4.0}

	\begin{tcolorbox}[colback=cyan,
		colframe=cyan,  
		arc=0pt,outer arc=0pt,
		valign=top, 
		halign=center,
		width=\textwidth]
		
		\includegraphics[width=.5cm]{../img/licence/cc_icon_white_x2.png}
		\includegraphics[width=.5cm]{../img/licence/attribution_icon_white_x2.png}
		
		\color{white}
		\bfseries Attribution 4.0 International (CC BY 4.0) \\
		\tiny \url{https://creativecommons.org/licenses/by/4.0/deed.fr}
		
	\end{tcolorbox}\vspace{-.5cm}
	\begin{tcolorbox}[colback=my-grey,
		colframe=my-grey,  
		center, arc=0pt,outer arc=0pt,
		valign=top, 
		halign=left,
		width=\textwidth]
		
		\tiny
		
		\begin{center}
			\bfseries\Large
			Vous êtes autorisé à :
		\end{center}
		
		\begin{minipage}{0.83\textwidth}
			\begin{itemize}
				\item[] \textbf{Partager} — copier, distribuer et communiquer le matériel par tous moyens et sous tous formats
				\item[] \textbf{Adapter} — remixer, transformer et créer à partir du matériel
				pour toute utilisation, y compris commerciale.
			\end{itemize}
		\end{minipage}
		\begin{minipage}{0.15\textwidth}
			\includegraphics[width=\textwidth]{../img/licence/FreeCulturalWorks_seal_x2.jpg}
		\end{minipage}
	
		
		\begin{center}
			\bfseries\Large
			Selon les conditions suivantes :
		\end{center}
		
		\begin{itemize}
			\item[] \textbf{Attribution} — Vous devez créditer l'Œuvre, intégrer un lien vers la licence et indiquer si des modifications ont été effectuées à l'Oeuvre. Vous devez indiquer ces informations par tous les moyens raisonnables, sans toutefois suggérer que l'Offrant vous soutient ou soutient la façon dont vous avez utilisé son Oeuvre. 
			\item[] \textbf{Pas de restrictions complémentaires} — Vous n'êtes pas autorisé à appliquer des conditions légales ou des mesures techniques qui restreindraient légalement autrui à utiliser l'Oeuvre dans les conditions décrites par la licence.
		\end{itemize}
		
	\end{tcolorbox}
	
%	\begin{center}
%		\bfseries Attribution 4.0 International (CC BY 4.0)
%		\url{https://creativecommons.org/licenses/by/4.0/deed.fr}
%	\end{center}

%	\tiny
%
%	Vous êtes autorisé à : 
%	\begin{itemize}
%		\item \textbf{Partager} — copier, distribuer et communiquer le matériel par tous moyens et sous tous formats
%		\item \textbf{Adapter} — remixer, transformer et créer à partir du matériel
%	\end{itemize}
%	
%	Selon les conditions suivantes : 
%	\begin{itemize}
%		\item \textbf{Attribution} — Vous devez créditer l'Œuvre, intégrer un lien vers la licence et indiquer si des modifications ont été effectuées à l'Oeuvre. Vous devez indiquer ces informations par tous les moyens raisonnables, sans toutefois suggérer que l'Offrant vous soutient ou soutient la façon dont vous avez utilisé son Oeuvre.
%		\item \textbf{Pas d'Utilisation Commerciale} — Vous n'êtes pas autorisé à faire un usage commercial de cette Oeuvre, tout ou partie du matériel la composant. 
%		\item \textbf{Pas de restrictions complémentaires} — Vous n'êtes pas autorisé à appliquer des conditions légales ou des mesures techniques qui restreindraient légalement autrui à utiliser l'Oeuvre dans les conditions décrites par la licence.
%	\end{itemize}

	\end{frame}
}

\settowidth{\leftmargini}{\usebeamertemplate{itemize item}}
\addtolength{\leftmargini}{\labelsep}

\AtBeginDocument{
	\newcolumntype{L}[2]{>{\vbox to #2\bgroup\vfill\flushleft}p{#1}<{\egroup}} 
	
	\begin{frame}[plain]
		\maketitle
	\end{frame}

	\insertlicence
}


% needs etoolbox; to break links after -
\appto\UrlBreaks{\do\-}


\makeatletter
\newcommand{\xRightarrow}[2][]{\ext@arrow 0359\Rightarrowfill@{#1}{#2}}
\makeatother


\usefonttheme{structurebold}
%\usefonttheme{professionalfonts}

\title[TALN : 09- Cohérence du discours]%
{Traitement automatique du langage naturel\\Chapitre 09 : Cohérence du discours} 

\changegraphpath{../img/coherence/}

\begin{document}
	
\begin{frame}
\frametitle{Traitement automatique du langage naturel}
\framesubtitle{Cohérence du discours : Introduction}

\begin{exampleblock}{Exemple d'un texte en français}
	\begin{center}
		\Large\bfseries
		L'hiver est un des quatre saisons de l'année. 
		Le rapport était bien exposé. 
		La pluie est une des caractéristiques de cette saison.
		Le sujet était la saison de l'hiver.
	\end{center}
\end{exampleblock}

\begin{itemize}
	\item Est-ce que vous pouvez comprendre de quoi s'agit-il ?
	\item S'il y a un problème, lequel ?
	\item Est-ce qu'on peut améliorer ce texte ?
\end{itemize}

\end{frame}

%\begin{frame}
%\frametitle{Traitement automatique du langage naturel}
%\framesubtitle{Sens des mots et désambigüisation lexicale : Un peu d'humour}
%
%\begin{center}
%	\vgraphpage{humour-parse.jpg}
%\end{center}
%
%\end{frame}

\begin{frame}
\frametitle{Traitement automatique du langage naturel}
\framesubtitle{Cohérence du discours : Plan}

\begin{multicols}{2}
%	\small
\tableofcontents
\end{multicols}
\end{frame}

%===================================================================================
\section{Relations de cohérence}
%===================================================================================

\begin{frame}
	\frametitle{Cohérence du discours}
	\framesubtitle{Relations de cohérence}
	
	Comment détecter la cohérence d'un discours ?
	\begin{itemize}
		\item \optword{Cohérence basée sur les relations} : les phrases d'un discours sont liées entre elles par des relations. 
		\begin{itemize}
			\item Rhetorical Structure Theory (RST) 
			\item Penn Discourse TreeBank (PDTB)
		\end{itemize}
		\item \optword{Cohérence basée sur l'entité} : les phrases d'un discours forment une seule entité ; elles discutent le même sujet
		\begin{itemize}
			\item Centering Theory 
		\end{itemize}
		\item \optword{Cohésion lexicale} : les phrases proches se partagent quelques mots ou sens 
		\begin{itemize}
			\item Similarité
		\end{itemize}
	\end{itemize}
	
\end{frame}

\subsection{Rhetorical Structure Theory (RST)}

\begin{frame}
\frametitle{Cohérence du discours : Relations de cohérence}
\framesubtitle{Rhetorical Structure Theory (RST)}
	
	\begin{itemize}
		\item Structure  
		\begin{itemize}
			\item \optword{Noyau (N)} :  unité indépendante, peut être interprétée indépendamment des autres unités du texte
			\item \optword{Satellite (S)} :  unité dépendante, ne peut être interprétée qu'en utilisant le noyau
		\end{itemize}
		\item Acteurs
		\begin{itemize}
			\item \optword{Lecteur (L)} :  qui a lit le script
			\item \optword{Scripteur (Sc)} :  qui a rédigé le script
		\end{itemize}
		\item Définition d'une relation
		\begin{itemize}
			\item Des restrictions sur le Noyau 
			\item Des restrictions sur le Satellite 
			\item Des restrictions sur la combinaison du Noyau et du Satellite 
			\item L'effet produit.
		\end{itemize}
	\end{itemize}
	
\end{frame}

\begin{frame}
	\frametitle{Cohérence du discours : Relations de cohérence}
	\framesubtitle{Rhetorical Structure Theory (RST) : Relations RST (Élaboration)}
	
	\begin{itemize}
		\item \optword{Élaboration} :  S donne des informations additionnelles sur la situation présentée dans N
		\begin{itemize}
			\item \textbf{Contraintes sur N + S :} S présente des informations supplémentaires vis-à-vis de la situation ou de quelqu'aspect du thème présenté(e) dans N 
			\item \textbf{Effet :}  L reconnaît la situation présentée dans S comme fournissant des informations supplémentaires au sujet de N. 
			\item \textbf{Lieu de l'effet :} N + S
			\item Ex. \expword{[\textsubscript{N} L'examen est facile.] [\textsubscript{S} Il ne prend qu'une heur.]}
		\end{itemize}
	\end{itemize}
	
\end{frame}

\begin{frame}
	\frametitle{Cohérence du discours : Relations de cohérence}
	\framesubtitle{Rhetorical Structure Theory (RST) : Relations RST (Évidence)}
	
	\begin{itemize}
		\item \optword{Évidence} :  S donne des informations additionnelles sur la situation présentée dans N dans le but de convincre L à accepter les informations de N
		\begin{itemize}
			\item \textbf{Contraintes sur N :} L pourra ne pas croire en N à un degré suffisant pour Sc
			\item \textbf{Contraintes sur S :} L croit S ou le trouve crédible
			\item \textbf{Contraintes sur N + S :} La compréhension de S par L augmente sa croyance en N
			\item \textbf{Effet :}  La croyance de L en N est effectivement augmentée
			\item \textbf{Lieu de l'effet :} N
			\item Ex. \expword{[\textsubscript{N} Kevin doit être ici.] [\textsubscript{S} Sa voiture est garée à l'extérieur.]}
		\end{itemize}
	\end{itemize}
	
\end{frame}

\begin{frame}
	\frametitle{Cohérence du discours : Relations de cohérence}
	\framesubtitle{Rhetorical Structure Theory (RST) : Relations RST (Contraste)}
	
	\begin{itemize}
		\item \optword{Contraste} :  Une relation d'opposition entre deux noyaux
		\begin{itemize}
			\item \textbf{Contraintes sur N :} multi-noyaux
			\item \textbf{Contraintes sur N + N :} pas plus de deux noyaux ; les situations présentées dans ces deux noyaux sont (a) comprises comme les mêmes à beaucoup d'égards, (b) comprises comme différant par quelques aspects et (c) comparées par rapport à l'une ou plus d'une de ces différences
			\item \textbf{Effet :}  L reconnaît la comparabilité et les différences fournies par la comparaison effectuée
			\item \textbf{Lieu de l'effet :} noyaux multiples
			\item Ex. \expword{[\textsubscript{N} Il a voté "Non" à la nouvelle constitution.] [\textsubscript{N} Son frère a voté "Oui".]}
		\end{itemize}
	\end{itemize}
	
\end{frame}

\begin{frame}
	\frametitle{Cohérence du discours : Relations de cohérence}
	\framesubtitle{Rhetorical Structure Theory (RST) : Relations RST (Antithèse)}
	
	\begin{itemize}
		\item \optword{Antithèse} :  Une relation d'opposition entre N et S dans le but que L ait une attitude positive envers N
		\begin{itemize}
			\item \textbf{Contraintes sur N :} Sc a une attitude positive par rapport à la situation présentée dans N
			\item \textbf{Contraintes sur N + N :} es situations présentées dans N et S sont en opposition (cf. CONTRASTE)
			\item \textbf{Effet :}  L'attitude positive de L vis-à-vis de N est augmentée
			\item \textbf{Lieu de l'effet :} N
			\item Ex. \expword{[\textsubscript{N} J'ai défendu la République jeune;] [\textsubscript{S} Je ne l'abandonnerai pas maintenant que je suis vieux.]}
		\end{itemize}
	\end{itemize}
	
\end{frame}

\begin{frame}
	\frametitle{Cohérence du discours : Relations de cohérence}
	\framesubtitle{Rhetorical Structure Theory (RST) : Relations thématiques vs. présentationnelles}
	
	\rowcolors{2}{lightblue}{lightyellow}\footnotesize
	\begin{tabular}{p{.4\textwidth}lp{.4\textwidth}}
		\rowcolor{darkblue}
		\bfseries\textcolor{white}{thématique} && \bfseries\textcolor{white}{présentationnelle}\\
		
		Élaboration
		
		Circonstance
		
		Problème-Solution
		
		Cause intentionnelle
		
		Résultat intentionnel
		
		Cause non-intentionnelle
		
		Résultat non-intentionnel
		
		But
		
		Condition
		
		Interprétation
		
		Évaluation
		
		Reformulation
		
		Résumé
		
		Séquence
		
		Contraste
		
		&&
		Motivation
		
		Antithèse
		
		Arrière-plan
		
		Facilitation
		
		Évidence (« indice »)
		
		Justification
		
		Concession\\
	\end{tabular}
	
\end{frame}

\begin{frame}
	\frametitle{Cohérence du discours : Relations de cohérence}
	\framesubtitle{Rhetorical Structure Theory (RST) : Relations sémantiques vs. pragmatiques vs. textuelles}
	
	\rowcolors{2}{lightblue}{lightyellow}\footnotesize
	\begin{tabular}{p{.3\textwidth}lp{.25\textwidth}lp{.2\textwidth}}
		\rowcolor{darkblue}
		\bfseries\textcolor{white}{Sémantique} && \bfseries\textcolor{white}{Pragmatique} && \bfseries\textcolor{white}{Textuelle} \\
		
		Élaboration
		
		Circonstance
		
		Problème-Solution
		
		Cause intentionnelle
		
		Résultat intentionnel
		
		Cause non-intentionnelle
		
		Résultat non-intentionnel
		
		But
		
		Condition
		
		Interprétation
		
		Évaluation
		
		Séquence
		
		Contraste
		
		&&
		
		Motivation
		
		Antithèse
		
		Facilitation
		
		Évidence (« indice »)
		
		Justification
		
		Concession
		
		&&
		
		Arrière-plan
		
		Reformulation
		
		Résumé\\
	\end{tabular}
	
\end{frame}

\subsection{Penn Discourse TreeBank (PDTB)}

\begin{frame}
	\frametitle{Cohérence du discours : Relations de cohérence}
	\framesubtitle{Penn Discourse TreeBank (PDTB)}
	
	\begin{itemize}
		\item Annotation des connectives de discours
		\begin{itemize}
			\item \optword{Conjonctions de subordination} :  
			temporelle (Ex., \expword{'when', 'as soon as'}), 
			causale (Ex., \expword{'because'}), 
			concessive (Ex., \expword{'although', 'even though'}), 
			objectif (Ex.,\expword{'so that', 'in order that'}) et 
			conditionnelle (Ex., \expword{'if', 'unless'}).
			
			\item \optword{Conjonctions de coordination} : \expword{'and', 'but', 'or'}
			
			\item \optword{Conjonctives adverbiales} : des adverbes qui expriment une relation de discours entre des évenements ou des états. Ex., \expword{'however', 'therefore', 'then', etc.}
			Des syntagmes prépositionnels sont inclus aussi dans cette classe. Ex. \expword{'as a result',
			'in addition', 'in fact', etc. }
			
			\item \optword{Connectives implicites} :  identifiées entre deux phrases adjacentes qui ne sont pas reliées par des connectives explicites.
		\end{itemize}
		\item Annotation des arguments 
		\begin{itemize}
			\item ARG2 : c'est la clause où la connective est liée (cas explicite)
			\item ARG1 : l'autre clause
		\end{itemize}
	\end{itemize}
	
\end{frame}

\begin{frame}
	\frametitle{Cohérence du discours : Relations de cohérence}
	\framesubtitle{Penn Discourse TreeBank (PDTB) : Sens des connectives}
	
	\begin{table}
	\rowcolors{2}{lightblue}{lightyellow}\tiny\bfseries
	\begin{tabular}{p{.1\textwidth}lp{.8\textwidth}}
		\rowcolor{darkblue}
		\bfseries\textcolor{white}{Connective} && \bfseries\textcolor{white}{Sens}\\
		
		after && succession (523), succession-reason (50), other (4) \\
		since && reason (94), succession (78), succession-reason (10), other (2) \\
		when && Synchrony (477), succession (157), general (100), succession-reason (65), Synchrony-general (50),
		Synchrony-reason (39), hypothetical (11), implicit assertion (11), Synchrony-hypothetical (10), other
		(69) \\
		while && juxtaposition (182), Synchrony (154), Contrast (120), expectation (79), opposition (78), Conjunction
		(39), Synchrony-juxtaposition (26), Synchrony-Conjunction (21), Synchrony-Contrast(22), COMPARISON (18), Synchrony-opposition (11), other (31) \\
		meanwhile && Synchrony-Conjunction (92), Synchrony (26), Conjunction (25), Synchrony-juxtaposition (15),
		other(35)\\
		but && Contrast (1609), juxtaposition (636), contra-expectation (494), COMPARISTON (260), opposition
		(174), Conjunction (63), Conjunction-Pragmatic contrast (14), Pragmatic-contrast (14), other (32)
		however Contrast (254), juxtaposition (89), contra-expectation (70), COMPARISON (49), opposition (31),
		other (12)\\
		although && expectation (132), Contrast (114) juxtaposition (34), contra-expectation (21), COMPARISON (16),
		opposition (9), other (2)\\
		and && Conjunction (2543), List (210), result-Conjunction (138), result (38), precedence-Conjunction (30),
		juxtaposition (11), other(30)\\
		if && hypothetical (682), general (175), unreal present (122), factual present (73), unreal past (53), expectation (34), implicit assertion (29), relevance (20), other (31)\\
	\end{tabular}
	\caption{Quelques connectives et des statistiques sur leurs sens \cite{2008-prasad-al}}
	\end{table}
	
\end{frame}

\begin{frame}
	\frametitle{Cohérence du discours : Relations de cohérence}
	\framesubtitle{Penn Discourse TreeBank (PDTB) : Sens des connectives}
	
	\begin{figure}
		\vgraphpage[.7\textheight]{pdtb-sens_.pdf}
		\caption{Hierarchie des sens dans PDTB \cite{2008-prasad-al}}
	\end{figure}
	
\end{frame}

%===================================================================================
\section{Analyse de discours}
%===================================================================================

\begin{frame}
	\frametitle{Cohérence du discours}
	\framesubtitle{Analyse de discours}
	
	\begin{itemize}
		\item 
	\end{itemize}
	
\end{frame}

\subsection{Segmentation des unités pour RST}

\begin{frame}
	\frametitle{Cohérence du discours : Analyse de discours}
	\framesubtitle{Segmentation des unités pour RST}
	
	\begin{itemize}
		\item 
	\end{itemize}
	
\end{frame}

\subsection{Analyse RST}

\begin{frame}
	\frametitle{Cohérence du discours : Analyse de discours}
	\framesubtitle{Analyse RST}
	
	\begin{itemize}
		\item 
	\end{itemize}
	
\end{frame}


\subsection{Analyse RST}

\begin{frame}
	\frametitle{Cohérence du discours : Analyse de discours}
	\framesubtitle{Analyse RST}
	
	\begin{itemize}
		\item 
	\end{itemize}
	
\end{frame}

\subsection{Analyse PDTB}

\begin{frame}
	\frametitle{Cohérence du discours : Analyse de discours}
	\framesubtitle{Analyse PDTB}
	
	\begin{itemize}
		\item 
	\end{itemize}
	
\end{frame}


\insertbibliography{TALN08}{*}

\end{document}

